\documentclass[letterpaper,10pt]{article}
\usepackage[margin=0.72in]{geometry}
\usepackage[T1]{fontenc}
\usepackage{lmodern}
\usepackage{booktabs}
\usepackage{longtable}
\usepackage{array}
\usepackage{float}
\usepackage{graphicx}
\usepackage{url}
\setlength{\parindent}{0pt}
\setlength{\parskip}{0.45em}
\renewcommand{\arraystretch}{1.08}

\title{Technical Appendix for Anonymous Review}
\author{}
\date{}

\begin{document}
\maketitle
\small

\section{Resource and Construction Details}

CliniCause provides two dataset-specific, estimator-ready causal-analysis
resources: MIMIC-III (26,845 analysis records) and PhysioNet 2012 (7,993).
The analysis unit is a source time-series record or ICU stay, rather than a
unique person.  At each source boundary, accepted identifiers are normalized to
the canonical \texttt{ts\_id} key; long irregular events retain minute,
variable, and value, while a record-level table carries in-hospital mortality
and available covariates.  The two resources share an interface but are never
pooled: their variable vocabularies, proxy ontologies, measurement processes,
and DAG assumptions remain source specific.

\begin{table}[H]
\centering\small
\caption{Admitted original-cohort exposure sets and source-specific semantics.}
\label{tab:supp-resources}
\begin{tabular}{p{0.19\linewidth}p{0.35\linewidth}p{0.37\linewidth}}
\toprule
Resource & Admitted exposures & Source-specific proxy boundary\\
\midrule
MIMIC-III & Cardiac strain; inflammation/sepsis; hepatic/coagulation dysfunction; renal dysfunction; global severity; metabolic derangement; neurologic dysfunction; respiratory failure; shock. & Hepatic and coagulation evidence are represented by a combined proxy.  Chronic-burden fields are retained as baseline constructs rather than admitted exposures.\\
PhysioNet 2012 & Cardiac injury/strain; coagulation/hematologic dysfunction; global severity; hepatic dysfunction; inflammation/sepsis burden; metabolic derangement; neurologic dysfunction; renal dysfunction; respiratory failure; shock. & Hepatic and coagulation/hematologic evidence are represented separately.  Chronic baseline constructs remain outside the admitted exposure set.\\
\bottomrule
\end{tabular}
\end{table}

At design time, structured LLM proposals covered proxy ontologies, rule
families, missingness considerations, and dataset-specific DAGs.  Project
selection encoded chosen proposals as deterministic proxy-rule and DAG source
artifacts.  The LLM did not receive runtime patient rows and was not a
patient-level estimator.  At construction time, each source is preprocessed in
its own lane; deterministic rule labels and four temporal models (STraTS, GRU,
GRU-D, and TCN) produce aligned annotation exports.  Each export is normalized
to probability and binary fields; one deterministic rule source and four model
sources are then combined by a fixed deterministic aggregation.  The final
resource joins the aggregate exposure layer to outcome, observed covariates,
DAG/adjustment metadata, and provenance records.

Validation is layered: schema handoffs check required columns, order, value
ranges, binary/probability consistency, and expected output sets; cohort
handoffs reject conflicting IDs, missing/extra records, and silent shrinkage;
prediction exports and aggregation enforce complete aligned voters; and
provenance artifacts bind available hashes, manifests, configurations, and
receipts.  Archived-result evidence establishes checked aggregate outputs, but
the exact historical producing revisions, configurations, and
checkpoint-to-export lineage remain incomplete.  Current repaired contracts are
documented separately and are not claimed to have produced the archived results.

\section{Complete Predictive Evaluation}

The predictive task is multi-label estimation of each dataset's deterministic
proxy vector.  Test loss is sample-mean multi-label binary cross-entropy.
AUROC, AUPRC, and minRP are calculated per target then macro-averaged after
excluding target columns with a single observed evaluation class; minRP is the
maximum over thresholds of the smaller of precision and recall.  The table
reports selected archived held-out point metrics.  It does not supply repeated
run uncertainty, significance testing, split-ID lineage, or a verified
checkpoint-to-export mapping.

\begin{table}[H] \centering \small
\caption{Complete selected archived held-out predictive results. Loss is minimized; the other metrics are maximized. Degenerate target columns are excluded from macro AUROC, AUPRC, and minRP.}\label{tab:supp-predictive}
\begin{tabular}{llrrrr} \toprule Dataset & Model & Loss & AUROC & AUPRC & minRP\\ \midrule
MIMIC-III & STRATS & 0.348136 & 0.905411 & 0.869417 & 0.795399\\
 & GRU & 0.385757 & 0.881119 & 0.839559 & 0.770421\\
 & GRU-D & 0.404448 & 0.884277 & 0.841273 & 0.772533\\
 & TCN & 0.414456 & 0.866914 & 0.822642 & 0.757798\\
PhysioNet 2012 & STRATS & 0.340692 & 0.914761 & 0.877456 & 0.818228\\
 & GRU & 0.396819 & 0.914616 & 0.897493 & 0.831044\\
 & GRU-D & 0.330593 & 0.918478 & 0.905105 & 0.834958\\
 & TCN & 0.476513 & 0.899103 & 0.875091 & 0.811184\\
\bottomrule \end{tabular}
\end{table}


\section{Complete Original-Cohort Estimator Matrix}

\begin{longtable}{@{}p{0.84in}p{1.42in}rrrp{0.68in}@{}}
\caption{Complete original-cohort estimator summaries (CF-DML denotes CausalForestDML). Values are mean model-estimated CATE summaries over the analyzed samples; display precision is six decimal places.}\label{tab:supp-original-cate}\\
\toprule Dataset & Exposure & CF-DML & LinearDML & CausalPFN & Sign agreement\\
\midrule \endfirsthead
\toprule Dataset & Exposure & CF-DML & LinearDML & CausalPFN & Sign agreement\\
\midrule \endhead
MIMIC-III & Cardiac strain & +0.220356 & +0.187634 & +0.259327 & all three\\
MIMIC-III & Inflammation/sepsis & +0.161179 & +0.161334 & +0.148544 & all three\\
MIMIC-III & Hepatic/coag. dysfunction & +0.097598 & +0.083017 & +0.096846 & all three\\
MIMIC-III & Renal dysfunction & +0.091290 & +0.088595 & +0.086972 & all three\\
MIMIC-III & Global severity & +0.062101 & +0.080153 & +0.073050 & all three\\
MIMIC-III & Respiratory failure & +0.036133 & +0.039098 & +0.049766 & all three\\
MIMIC-III & Neurologic dysfunction & +0.034253 & +0.031550 & +0.034993 & all three\\
MIMIC-III & Shock & +0.020994 & +0.021246 & +0.031645 & all three\\
MIMIC-III & Metabolic derangement & +0.019743 & +0.017989 & +0.031323 & all three\\
PhysioNet 2012 & Renal dysfunction & +0.120027 & +0.090622 & +0.110669 & all three\\
PhysioNet 2012 & Cardiac injury/strain & +0.111831 & +0.122033 & +0.119413 & all three\\
PhysioNet 2012 & Global severity & +0.107988 & +0.106740 & +0.105039 & all three\\
PhysioNet 2012 & Hepatic dysfunction & +0.090527 & +0.060184 & +0.106843 & all three\\
PhysioNet 2012 & Neurologic dysfunction & +0.081636 & +0.075728 & +0.077049 & all three\\
PhysioNet 2012 & Metabolic derangement & +0.077704 & +0.079738 & +0.090734 & all three\\
PhysioNet 2012 & Inflammation/sepsis burden & +0.071800 & +0.070872 & +0.067480 & all three\\
PhysioNet 2012 & Respiratory failure & +0.064750 & +0.062660 & +0.062984 & all three\\
PhysioNet 2012 & Coag./heme dysfunction & +0.010794 & +0.004136 & +0.025428 & all three\\
PhysioNet 2012 & Shock & -0.013849 & -0.026944 & +0.004122 & disagreement\\
\bottomrule \end{longtable}


\section{Outcome-Downsampling Comparison}

All mortality-positive records are retained while mortality-negative records are
sampled under the implemented configuration.  This produces a different
analyzed population and outcome prevalence; therefore, the table uses direction
as a robustness comparison and does not treat downsampled magnitudes as
transportable estimates for the original populations.  The generator verifies
all 57 matched original/downsampled estimator rows and exactly two sign changes.

\begin{longtable}{@{}p{0.78in}p{1.18in}p{0.56in}rrlp{0.62in}@{}}
\caption{Original versus outcome-downsampled estimator summaries (CF-DML, Linear, and PFN denote CausalForestDML, LinearDML, and CausalPFN). Direction compares signs only; the changed population means magnitudes are not transported between conditions.}\label{tab:supp-downsampling}\\
\toprule Dataset & Exposure & Estimator & Original & Downsampled & Original sign & Direction\\
\midrule \endfirsthead
\toprule Dataset & Exposure & Estimator & Original & Downsampled & Original sign & Direction\\
\midrule \endhead
MIMIC-III & Cardiac strain & CF-DML & +0.220356 & +0.420725 & + & preserved\\
MIMIC-III & Cardiac strain & Linear & +0.187634 & +0.403937 & + & preserved\\
MIMIC-III & Cardiac strain & PFN & +0.259327 & +0.446536 & + & preserved\\
MIMIC-III & Global severity & CF-DML & +0.062101 & +0.154682 & + & preserved\\
MIMIC-III & Global severity & Linear & +0.080153 & +0.065034 & + & preserved\\
MIMIC-III & Global severity & PFN & +0.073050 & +0.190187 & + & preserved\\
MIMIC-III & Hepatic/coag. dysfunction & CF-DML & +0.097598 & +0.192393 & + & preserved\\
MIMIC-III & Hepatic/coag. dysfunction & Linear & +0.083017 & +0.182470 & + & preserved\\
MIMIC-III & Hepatic/coag. dysfunction & PFN & +0.096846 & +0.219766 & + & preserved\\
MIMIC-III & Inflammation/sepsis & CF-DML & +0.161179 & +0.407940 & + & preserved\\
MIMIC-III & Inflammation/sepsis & Linear & +0.161334 & +0.393109 & + & preserved\\
MIMIC-III & Inflammation/sepsis & PFN & +0.148544 & +0.419575 & + & preserved\\
MIMIC-III & Metabolic derangement & CF-DML & +0.019743 & +0.028995 & + & preserved\\
MIMIC-III & Metabolic derangement & Linear & +0.017989 & +0.026342 & + & preserved\\
MIMIC-III & Metabolic derangement & PFN & +0.031323 & +0.052237 & + & preserved\\
MIMIC-III & Neurologic dysfunction & CF-DML & +0.034253 & +0.074812 & + & preserved\\
MIMIC-III & Neurologic dysfunction & Linear & +0.031550 & +0.067685 & + & preserved\\
MIMIC-III & Neurologic dysfunction & PFN & +0.034993 & +0.076865 & + & preserved\\
MIMIC-III & Renal dysfunction & CF-DML & +0.091290 & +0.180319 & + & preserved\\
MIMIC-III & Renal dysfunction & Linear & +0.088595 & +0.180638 & + & preserved\\
MIMIC-III & Renal dysfunction & PFN & +0.086972 & +0.203699 & + & preserved\\
MIMIC-III & Respiratory failure & CF-DML & +0.036133 & +0.083960 & + & preserved\\
MIMIC-III & Respiratory failure & Linear & +0.039098 & +0.090434 & + & preserved\\
MIMIC-III & Respiratory failure & PFN & +0.049766 & +0.129317 & + & preserved\\
MIMIC-III & Shock & CF-DML & +0.020994 & +0.021534 & + & preserved\\
MIMIC-III & Shock & Linear & +0.021246 & +0.013945 & + & preserved\\
MIMIC-III & Shock & PFN & +0.031645 & +0.036761 & + & preserved\\
PhysioNet 2012 & Cardiac injury/strain & CF-DML & +0.111831 & +0.225488 & + & preserved\\
PhysioNet 2012 & Cardiac injury/strain & Linear & +0.122033 & +0.243940 & + & preserved\\
PhysioNet 2012 & Cardiac injury/strain & PFN & +0.119413 & +0.240019 & + & preserved\\
PhysioNet 2012 & Coag./heme dysfunction & CF-DML & +0.010794 & +0.011005 & + & preserved\\
PhysioNet 2012 & Coag./heme dysfunction & Linear & +0.004136 & -0.017406 & + & changed\\
PhysioNet 2012 & Coag./heme dysfunction & PFN & +0.025428 & +0.055540 & + & preserved\\
PhysioNet 2012 & Global severity & CF-DML & +0.107988 & +0.244412 & + & preserved\\
PhysioNet 2012 & Global severity & Linear & +0.106740 & +0.206430 & + & preserved\\
PhysioNet 2012 & Global severity & PFN & +0.105039 & +0.314587 & + & preserved\\
PhysioNet 2012 & Hepatic dysfunction & CF-DML & +0.090527 & +0.214685 & + & preserved\\
PhysioNet 2012 & Hepatic dysfunction & Linear & +0.060184 & +0.147114 & + & preserved\\
PhysioNet 2012 & Hepatic dysfunction & PFN & +0.106843 & +0.216485 & + & preserved\\
PhysioNet 2012 & Inflammation/sepsis burden & CF-DML & +0.071800 & +0.147896 & + & preserved\\
PhysioNet 2012 & Inflammation/sepsis burden & Linear & +0.070872 & +0.141951 & + & preserved\\
PhysioNet 2012 & Inflammation/sepsis burden & PFN & +0.067480 & +0.150122 & + & preserved\\
PhysioNet 2012 & Metabolic derangement & CF-DML & +0.077704 & +0.187367 & + & preserved\\
PhysioNet 2012 & Metabolic derangement & Linear & +0.079738 & +0.192312 & + & preserved\\
PhysioNet 2012 & Metabolic derangement & PFN & +0.090734 & +0.228162 & + & preserved\\
PhysioNet 2012 & Neurologic dysfunction & CF-DML & +0.081636 & +0.163649 & + & preserved\\
PhysioNet 2012 & Neurologic dysfunction & Linear & +0.075728 & +0.138697 & + & preserved\\
PhysioNet 2012 & Neurologic dysfunction & PFN & +0.077049 & +0.182251 & + & preserved\\
PhysioNet 2012 & Renal dysfunction & CF-DML & +0.120027 & +0.226767 & + & preserved\\
PhysioNet 2012 & Renal dysfunction & Linear & +0.090622 & +0.192089 & + & preserved\\
PhysioNet 2012 & Renal dysfunction & PFN & +0.110669 & +0.224075 & + & preserved\\
PhysioNet 2012 & Respiratory failure & CF-DML & +0.064750 & +0.110970 & + & preserved\\
PhysioNet 2012 & Respiratory failure & Linear & +0.062660 & +0.095604 & + & preserved\\
PhysioNet 2012 & Respiratory failure & PFN & +0.062984 & +0.131635 & + & preserved\\
PhysioNet 2012 & Shock & CF-DML & -0.013849 & -0.017995 & - & preserved\\
PhysioNet 2012 & Shock & Linear & -0.026944 & -0.056043 & - & preserved\\
PhysioNet 2012 & Shock & PFN & +0.004122 & -0.041287 & + & changed\\
\bottomrule \end{longtable}


\section{Matching Availability and Support}

Matching is a descriptive, support-oriented, greedy one-to-one comparison over
the archived binary matching representation.  Successful rows report the
descriptive matched-pair outcome difference.  A failure indicates that no usable
binary matching columns remained after preprocessing; it is not an effect of
zero.

\begin{longtable}{@{}p{0.78in}p{1.10in}p{0.44in}p{0.48in}p{0.54in}p{0.48in}p{1.02in}@{}}
\caption{Original-cohort matching availability. Failures are retained as failures, not zero effects.}\label{tab:supp-matching}\\
\toprule Dataset & Exposure & Status & Pairs & Match rate & Final distance & Result / reason\\
\midrule \endfirsthead
\toprule Dataset & Exposure & Status & Pairs & Match rate & Final distance & Result / reason\\
\midrule \endhead
MIMIC-III & Cardiac strain & success & 4769 & 1.000 & 0 & +0.344726\\
MIMIC-III & Global severity & success & 7537 & 0.390 & 1 & +0.074831\\
MIMIC-III & Hepatic/coag. dysfunction & success & 6478 & 1.000 & 0 & +0.133374\\
MIMIC-III & Inflammation/sepsis & failure & -- & -- & -- & No binary matching columns after preprocessing\\
MIMIC-III & Metabolic derangement & success & 5061 & 0.506 & 0 & +0.036949\\
MIMIC-III & Neurologic dysfunction & success & 6275 & 0.539 & 0 & +0.026932\\
MIMIC-III & Renal dysfunction & success & 8915 & 0.915 & 0 & +0.133595\\
MIMIC-III & Respiratory failure & success & 10998 & 0.694 & 2 & +0.105656\\
MIMIC-III & Shock & success & 10380 & 0.717 & 1 & +0.038632\\
PhysioNet 2012 & Cardiac injury/strain & success & 2840 & 1.000 & 0 & +0.133803\\
PhysioNet 2012 & Coag./heme dysfunction & success & 1763 & 0.558 & 0 & +0.001702\\
PhysioNet 2012 & Global severity & success & 1782 & 0.287 & 1 & +0.095398\\
PhysioNet 2012 & Hepatic dysfunction & success & 1133 & 1.000 & 0 & +0.136805\\
PhysioNet 2012 & Inflammation/sepsis burden & failure & -- & -- & -- & No binary matching columns after preprocessing\\
PhysioNet 2012 & Metabolic derangement & success & 3378 & 0.757 & 2 & +0.073416\\
PhysioNet 2012 & Neurologic dysfunction & failure & -- & -- & -- & No binary matching columns after preprocessing\\
PhysioNet 2012 & Renal dysfunction & success & 1841 & 1.000 & 0 & +0.154264\\
PhysioNet 2012 & Respiratory failure & failure & -- & -- & -- & No binary matching columns after preprocessing\\
PhysioNet 2012 & Shock & success & 2433 & 0.438 & 1 & +0.010275\\
\bottomrule \end{longtable}


\section{Sensitivity and Permutation Coverage}

\begin{table}[H] \centering \small
\caption{Diagnostic coverage and provenance. Counts are generated from checked aggregate diagnostic records, not inferred from estimator names.}\label{tab:supp-diagnostic-coverage}
\begin{tabular}{p{0.22\linewidth}p{0.68\linewidth}} \toprule Coverage family & Checked status / boundary\\ \midrule
Estimator-native sensitivity & Archived DML records distinguish saved-training direct, recomputed/reconstructed, partial, and failed diagnostics; they are not collapsed.\\
CausalPFN sensitivity & No equivalent archived CausalPFN sensitivity stage; unavailable by archived-stage design.\\
Original-cohort sensitivity rows & CausalForestDML: FAILED=4; CausalForestDML: PARTIAL=15; LinearDML: FAILED=4; LinearDML: PARTIAL=15.\\
Treatment permutations & Available only for archived non-PFN estimator stages; disruption diagnostic, not identification proof.\\
Outcome permutations & Available only for archived non-PFN estimator stages; disruption diagnostic, not identification proof.\\
Original-cohort permutation rows & CausalForestDML: outcome=19; CausalForestDML: treatment=19; LinearDML: outcome=19; LinearDML: treatment=19.\\
\bottomrule \end{tabular}
\end{table}


\section{Provenance and Reproducibility Boundaries}

Checked archived artifacts support the resource counts, selected predictive
metrics, original and downsampled estimator summaries, matching availability,
and the recorded diagnostic status classes.  They do not establish an exact
historical clean-checkout rerun.  The known gaps include producing source
revisions, numbered causal configurations, predictive split IDs,
checkpoint-to-export linkage, raw/processed-data manifests, producing
environment/hardware capture, and CausalPFN package/version alignment.

Current source inspection documents stricter runtime validation gates for
schemas, canonical identifiers, voters, and artifact reuse.  Those current
contracts must be tested and release-validated separately; this appendix makes
no current test-pass claim.  Any eventual anonymous package must keep
source-data access restricted to authorized MIMIC-III and PhysioNet users,
include only licensing-permitted derived artifacts, and undergo a final identity
and metadata audit.  No patient-level values are included in this appendix.

\end{document}
